\documentclass[12pt,a4paper]{article}

\usepackage[T1]{fontenc}
\usepackage[utf8]{inputenc}
\usepackage[shorthands=off,brazilian]{babel}
\usepackage[a4paper,margin=2.5cm]{geometry}
\usepackage{booktabs}
\usepackage{tabularx}
\usepackage{longtable}
\usepackage{ragged2e}
\usepackage[hidelinks]{hyperref}
\usepackage[backend=biber,style=abnt,useprefix=true]{biblatex}

\addbibresource{referencias.bib}

\setlength{\parindent}{1.25cm}
\newcolumntype{Y}{>{\RaggedRight\arraybackslash}X}

\title{\textbf{Desigualdade geográfica no acesso a serviços de inteligência
artificial em nuvem:}\\[0.3em]
\large como a literatura mede defasagem de adoção e assimetria de acesso entre países}
\author{Pedro Dias Soares \and Eduardo Velloso Landeiro}
\date{Projeto em Ciência de Dados IV --- PUC Minas\\
Prof. Rodrigo Richard Gomes --- 2\textsuperscript{o} semestre de 2026\\[0.5em]
\today}

\begin{document}
\maketitle

\section{Tema}

Esta revisão trata da \textbf{desigualdade geográfica no acesso a serviços de
inteligência artificial em nuvem}, com foco em como a literatura mede a defasagem de
adoção e a assimetria de acesso entre países.

O recorte é deliberadamente metodológico. O objeto não é a desigualdade em si, já
amplamente descrita, mas o \emph{instrumento} com que ela é medida: que unidade de
análise se adota, de que fonte de dados se extrai a evidência e que método se aplica
para transformar acesso desigual em quantidade comparável entre países.

\section{Introdução}

Modelos de inteligência artificial generativa passaram, em poucos anos, de objeto de
pesquisa a insumo produtivo. Sistemas de saúde, universidades, empresas e órgãos
públicos constroem aplicações sobre modelos ofertados por um número reduzido de
empresas, concentradas em um único país, e disponibilizados por meio de plataformas de
computação em nuvem. Esse insumo não chega a todos os territórios ao mesmo tempo, pelo
mesmo preço nem com os mesmos recursos.

A literatura sobre desigualdade digital tem tradição consolidada. Desde a formulação
dos níveis do \emph{digital divide}, o campo documentou que a assimetria não se esgota
no acesso físico à rede, mas se estende à qualidade do equipamento, às habilidades de
uso e aos resultados obtidos \cite{vandeursen2018,lutz2019}. Ocorre que essa tradição
foi construída sobre um objeto --- o indivíduo conectado --- que não coincide com a camada
em que a assimetria de inteligência artificial parece operar hoje: a disponibilidade
de capacidade computacional de fronteira em um território.

Esta revisão examina dez trabalhos que, de ângulos distintos, enfrentam o problema de
medir acesso desigual à tecnologia. O objetivo não é resumi-los, mas confrontar seus
instrumentos de medição, identificar o que cada um consegue e o que não consegue
observar, e localizar a lacuna que resulta desse confronto.

\section{Estratégia de busca bibliográfica}

\subsection{Procedimento}

A busca seguiu o procedimento de \textcite{santos2007}: seleção de descritores por
componente da pergunta, combinação com operadores booleanos --- \texttt{OR} aditivo
dentro do componente, \texttt{AND} restritivo entre componentes --- e composição da
estratégia final pelo encadeamento dos componentes.

Uma adaptação foi necessária e é declarada aqui. Os vocabulários controlados citados
por \citeauthor{santos2007} --- MeSH, DeCS e EMTREE --- pertencem à área da saúde e não
se aplicam. Os equivalentes em Computação são o \emph{ACM Computing Classification
System} e o \emph{IEEE Thesaurus}, adotados como referência terminológica.

Registra-se também o tratamento dado ao componente de comparação. Em revisão
exploratória sobre tema emergente, incluí-lo na expressão de busca reduz
excessivamente a recuperação. O componente foi mantido \emph{implícito} na busca e
explícito no desenho da análise, conforme a própria ressalva de
\citeauthor{santos2007} de que a comparação pode ser nula.

\subsection{Bases consultadas}

A recuperação foi conduzida principalmente pela API da \textbf{OpenAlex}, base
bibliográfica aberta que indexa periódicos, atas de conferência e repositórios de
preprints, incluindo títulos SciELO. A verificação de existência das referências
recuperadas usou as APIs da \textbf{Crossref}, da \textbf{DataCite} e do
\textbf{arXiv}. Buscas em português foram executadas com filtro de idioma.

Duas limitações de acesso devem ser declaradas. Primeiro, \textbf{Scopus} e
\textbf{Web of Science} exigem assinatura institucional e não foram consultadas
diretamente; sua cobertura foi aproximada pela OpenAlex, que indexa em larga medida os
mesmos periódicos, mas a substituição não é equivalente e a busca deve ser repetida
com acesso institucional. Segundo, a busca direta no \textbf{SciELO} não respondeu a
partir do ambiente utilizado, e a cobertura em português foi suprida pelo filtro de
idioma da OpenAlex.

\subsection{Expressões de busca}

Os descritores foram organizados em três componentes:

\begin{quote}
\footnotesize
\texttt{P = ("digital divide" OR "digital inequality" OR "technology adoption"\\
\phantom{P = (}OR "cloud computing" OR "generative AI" OR "large language model")}

\vspace{0.4em}
\texttt{E = ("digital sovereignty" OR "AI sovereignty" OR "data residency"\\
\phantom{E = (}OR "export control" OR "compute governance" OR "global south")}

\vspace{0.4em}
\texttt{O = ("adoption lag" OR "diffusion" OR "access gap" OR "compute divide"\\
\phantom{O = (}OR "uneven development" OR "availability")}
\end{quote}

\noindent Expressão final: \texttt{(P) AND (E) AND (O)}, com exclusão de
\texttt{("federated learning" OR "edge computing" OR "network topology")}.

\noindent Em português:
\begin{quote}
\footnotesize
\texttt{("soberania digital" OR "desigualdade digital"\\
\phantom{(}OR "dependência tecnológica")\\
AND ("inteligência artificial" OR "computação em nuvem")}
\end{quote}

\noindent Filtros aplicados: publicações a partir de 2018; artigos, atas de
conferência, preprints e relatórios institucionais.

\subsection{Registro da busca}

Foram executadas 17 expressões de recuperação e 13 consultas de verificação, todas em
11 de setembro de 2026. O registro completo --- base, data, expressão exata, número de
retornos, número após triagem e observações --- está no arquivo \texttt{log-busca.csv}
que acompanha este trabalho. A Tabela~\ref{tab:busca} resume as expressões que
efetivamente produziram trabalhos selecionados.

\begin{table}[htbp]
\centering
\footnotesize
\caption{Expressões de busca que resultaram em seleção}
\label{tab:busca}
\begin{tabularx}{\textwidth}{Y r r}
\toprule
\textbf{Expressão (campo título e resumo)} & \textbf{Retornos} & \textbf{Selec.} \\
\midrule
\texttt{"compute divide"} & 45 & 1 \\
\texttt{"AI divide" OR "artificial intelligence divide"} & 268 & 2 \\
\texttt{"cross-country" AND "technology adoption" AND diffusion} & 22 & 1 \\
\texttt{"digital divide" AND ("artificial intelligence" OR "machine learning")} & 4.537 & 1 \\
\texttt{compute governance AI export controls semiconductors} & 1.638 & 1 \\
\texttt{"AI sovereignty" OR "sovereign AI" OR "digital sovereignty"} & 846 & 1 \\
\texttt{("large language model" OR "generative AI") AND (inequality OR "global south")} & 2.303 & 1 \\
\texttt{"Internet Archive" OR "Wayback Machine" AND (longitudinal OR measurement)} & 61 & 1 \\
\texttt{"soberania digital"} (idioma: português) & 98 & 1 \\
\midrule
\multicolumn{2}{l}{\textbf{Total selecionado}} & \textbf{10} \\
\bottomrule
\end{tabularx}
\end{table}

\subsection{Critérios de seleção}

Foram incluídos trabalhos acadêmicos --- artigos em periódico, atas de conferência,
preprints em repositório reconhecido --- diretamente relacionados à medição de acesso
desigual à tecnologia. Foram excluídos trabalhos sobre aprendizado federado,
computação de borda e topologia de rede, recorrentes nas expressões de busca mas
alheios ao objeto, e trabalhos sobre aplicações setoriais de IA sem tratamento da
questão distributiva.

Dos dez trabalhos selecionados, \textbf{oito foram publicados a partir de 2021},
satisfazendo o critério de que a maioria tenha no máximo cinco anos. Os dois
anteriores --- \textcite{lutz2019} e \textcite{ahmed2020} --- foram mantidos por
função específica: o primeiro estrutura o campo de desigualdades digitais e o segundo
é o trabalho que introduz a noção de \emph{compute divide}.

\begin{sloppypar}
Todas as referências tiveram existência verificada contra o registro da agência
mantenedora do identificador. O procedimento e suas ocorrências estão em
\texttt{verificacao-referencias.md}.
\end{sloppypar}

\subsection{Limitação da leitura}

A caracterização dos trabalhos registrada nas Tabelas~\ref{tab:compA} a~\ref{tab:compC} foi
construída a partir dos resumos e metadados obtidos nas APIs consultadas. As colunas
\emph{Principais resultados} e \emph{Limitações observadas} refletem, portanto, o que
os autores declaram em resumo, e não o que a leitura integral permitiria apurar. Onde
a limitação foi inferida do desenho do estudo, e não declarada pelos autores, isso
está assinalado.

\section{Trabalhos selecionados}

Os dez trabalhos distribuem-se em três correntes e uma vertente metodológica.

\paragraph{Corrente 1 --- Desigualdade digital.}
\begin{enumerate}
  \item \textcite{lutz2019} --- \emph{Human Behavior and Emerging Technologies}
  \item \textcite{mannuru2023} --- \emph{Information Development}
  \item \textcite{wang2024} --- \emph{New Media \& Society}
\end{enumerate}

\paragraph{Corrente 2 --- Difusão de tecnologia e defasagem de adoção.}
\begin{enumerate}\setcounter{enumi}{3}
  \item \textcite{knez2023} --- \emph{Journal of Evolutionary Economics}
  \item \textcite{mcelheran2024} --- \emph{Journal of Economics \& Management Strategy}
\end{enumerate}

\paragraph{Corrente 3 --- Governança de computação e soberania digital.}
\begin{enumerate}\setcounter{enumi}{5}
  \item \textcite{ahmed2020} --- arXiv, cs.CY
  \item \textcite{sastry2024} --- arXiv, cs.CY
  \item \textcite{mugge2024} --- \emph{Journal of European Public Policy}
  \item \textcite{furtado2024} --- \emph{Liinc em Revista}
\end{enumerate}

\paragraph{Vertente metodológica --- Arquivos da web como série temporal.}
\begin{enumerate}\setcounter{enumi}{9}
  \item \textcite{amos2021} --- \emph{Proceedings of the Web Conference} (ACM)
\end{enumerate}

\clearpage
\section{Tabela comparativa}

A comparação está distribuída em três tabelas por razão de largura de página. Juntas,
elas contêm as sete colunas exigidas pelo enunciado: a Tabela~\ref{tab:compA} traz
autoria e ano, título e problema investigado; a Tabela~\ref{tab:compB} traz método e
base de dados; a Tabela~\ref{tab:compC} traz principais resultados e limitações
observadas. A coluna de autoria e ano se repete nas três como chave de leitura, e a
ordem dos trabalhos é a mesma em todas.

\footnotesize
\begin{longtable}{@{}p{2.4cm} p{6.0cm} p{6.0cm}@{}}
\caption{Comparação dos trabalhos selecionados --- identificação e problema}\label{tab:compA}\\
\toprule
\textbf{Autor(es) / Ano} & \textbf{Título} & \textbf{Problema investigado} \\
\midrule
\endfirsthead
\multicolumn{3}{@{}l}{\footnotesize\itshape continuação da Tabela \thetable}\\
\toprule
\textbf{Autor(es) / Ano} & \textbf{Título} & \textbf{Problema investigado} \\
\midrule
\endhead
\midrule
\multicolumn{3}{r@{}}{\footnotesize\itshape continua na próxima página}\\
\endfoot
\bottomrule
\endlastfoot
Lutz (2019) & Digital inequalities in the age of artificial intelligence and big data & Como o campo de desigualdades digitais deve se reposicionar diante de \emph{big data} e inteligência artificial \\
\addlinespace
Amos, Acar, Lucherini \emph{et al.} (2021) & Privacy policies over time: curation and analysis of a million-document dataset & A ausência de base longitudinal impede analisar como documentos publicados na web evoluem ao longo do tempo \\
\addlinespace
Knez (2023) & Technology diffusion and uneven development & Modelos de difusão não consideram que o custo relativo de implementar uma tecnologia difere entre países \\
\addlinespace
Mannuru, Shahriar, Teel \emph{et al.} (2023) & Artificial intelligence in developing countries & Qual o impacto da inteligência artificial generativa sobre países em desenvolvimento \\
\addlinespace
Wang, Boerman, Kroon \emph{et al.} (2024) & The artificial intelligence divide: who is the most vulnerable? & Quais grupos são mais vulneráveis em um ambiente digital moldado por inteligência artificial \\
\addlinespace
McElheran, Li, Brynjolfsson \emph{et al.} (2024) & AI adoption in America: who, what, and where & Quem adota inteligência artificial nos Estados Unidos, quais tecnologias e em que locais \\
\addlinespace
Ahmed e Wahed (2020) & The de-democratization of AI: deep learning and the compute divide & O acesso desigual a poder computacional concentra a produção de pesquisa em inteligência artificial \\
\addlinespace
Sastry, Heim, Belfield \emph{et al.} (2024) & Computing power and the governance of artificial intelligence & De que modo o poder computacional pode operar como instrumento de governança de inteligência artificial \\
\addlinespace
Mügge (2024) & EU AI sovereignty: for whom, to what end, and to whose benefit? & A noção de soberania de inteligência artificial da União Europeia serve a objetivos diversos e conflitantes \\
\addlinespace
Furtado e Cunha (2024) & Inteligência artificial, data centers e colonialismo digital & A dimensão ecológica está ausente do debate sobre soberania digital \\
\addlinespace
\end{longtable}

\begin{longtable}{@{}p{2.4cm} p{6.0cm} p{6.0cm}@{}}
\caption{Comparação dos trabalhos selecionados --- método e base de dados}\label{tab:compB}\\
\toprule
\textbf{Autor(es) / Ano} & \textbf{Método ou abordagem} & \textbf{Base de dados utilizada} \\
\midrule
\endfirsthead
\multicolumn{3}{@{}l}{\footnotesize\itshape continuação da Tabela \thetable}\\
\toprule
\textbf{Autor(es) / Ano} & \textbf{Método ou abordagem} & \textbf{Base de dados utilizada} \\
\midrule
\endhead
\midrule
\multicolumn{3}{r@{}}{\footnotesize\itshape continua na próxima página}\\
\endfoot
\bottomrule
\endlastfoot
Lutz (2019) & Revisão de literatura; ensaio conceitual & Estudo conceitual, sem dados primários \\
\addlinespace
Amos, Acar, Lucherini \emph{et al.} (2021) & Construção de \emph{crawler} sobre a Wayback Machine, com etapas de validação e controle de qualidade & 1.071.488 políticas de privacidade em inglês, de mais de 130.000 sites, cobrindo cerca de vinte anos, extraídas do Internet Archive \\
\addlinespace
Knez (2023) & Modelo macroeconômico dinâmico de difusão, calibrado por aproximação não linear entre países e tecnologias & CHAT (\emph{Cross-country Historical Adoption of Technology}) e Penn World Table \\
\addlinespace
Mannuru, Shahriar, Teel \emph{et al.} (2023) & Revisão exploratória multiautoral & Estudo conceitual, sem dados primários \\
\addlinespace
Wang, Boerman, Kroon \emph{et al.} (2024) & Survey com análise de classes latentes & 1.088 cidadãos neerlandeses com 16 anos ou mais \\
\addlinespace
McElheran, Li, Brynjolfsson \emph{et al.} (2024) & Análise empírica descritiva sobre levantamento censitário de firmas & \emph{Annual Business Survey} de 2018, com cerca de 850.000 firmas dos Estados Unidos \\
\addlinespace
Ahmed e Wahed (2020) & Análise bibliométrica com estimador de controle sintético generalizado & 171.394 artigos de 57 conferências de ciência da computação \\
\addlinespace
Sastry, Heim, Belfield \emph{et al.} (2024) & Análise conceitual de instrumentos de política & Estudo conceitual, sem dados primários \\
\addlinespace
Mügge (2024) & Análise conceitual com mapeamento documental & Corpus documental: estratégia da Comissão Europeia desde 2018 e negociações do AI Act \\
\addlinespace
Furtado e Cunha (2024) & Revisão bibliográfica e pesquisa documental & Corpus documental sobre a expansão de data centers no Brasil entre 2013 e 2023 \\
\addlinespace
\end{longtable}

\begin{longtable}{@{}p{2.4cm} p{6.0cm} p{6.0cm}@{}}
\caption{Comparação dos trabalhos selecionados --- resultados e limitações}\label{tab:compC}\\
\toprule
\textbf{Autor(es) / Ano} & \textbf{Principais resultados} & \textbf{Limitações observadas} \\
\midrule
\endfirsthead
\multicolumn{3}{@{}l}{\footnotesize\itshape continuação da Tabela \thetable}\\
\toprule
\textbf{Autor(es) / Ano} & \textbf{Principais resultados} & \textbf{Limitações observadas} \\
\midrule
\endhead
\midrule
\multicolumn{3}{r@{}}{\footnotesize\itshape continua na próxima página}\\
\endfoot
\bottomrule
\endlastfoot
Lutz (2019) & Organiza o campo em três níveis --- acesso, habilidades de uso e resultados --- e estende o primeiro nível a sistemas dotados de inteligência artificial & Revisão sem protocolo de busca declarado; a unidade de análise é o indivíduo, o que não permite inferência sobre territórios \emph{(inferida)} \\
\addlinespace
Amos, Acar, Lucherini \emph{et al.} (2021) & Os documentos dobraram de extensão e tornaram-se menos legíveis ao longo de vinte anos; falham de forma consistente em divulgar tecnologias de rastreamento & Restrito ao inglês; a cobertura do Internet Archive é desigual entre sites, o que condiciona a série reconstruída \\
\addlinespace
Knez (2023) & O modelo calibrado explica cerca de 50\% da variabilidade da densidade tecnológica; as curvas de adoção são específicas por país e dependem do nível de desenvolvimento e do nível salarial & Cerca de metade da variabilidade permanece não explicada; as tecnologias do CHAT são majoritariamente industriais e anteriores à computação em nuvem \emph{(inferida)} \\
\addlinespace
Mannuru, Shahriar, Teel \emph{et al.} (2023) & Os benefícios da inteligência artificial generativa não são igualmente acessíveis; o acesso limitado a tecnologias de ponta e a infraestrutura inadequada são as barreiras apontadas & Não apresenta medição empírica; não distingue barreira de infraestrutura de barreira de oferta \emph{(inferida)} \\
\addlinespace
Wang, Boerman, Kroon \emph{et al.} (2024) & Identifica cinco perfis de usuário; os mais vulneráveis são mais velhos, menos escolarizados e com menor habilidade de proteção de privacidade, o que evidencia uma \emph{AI divide} emergente & Restrito a um único país; competências medidas por autorrelato; nível individual de análise \\
\addlinespace
McElheran, Li, Brynjolfsson \emph{et al.} (2024) & Menos de 6\% das firmas usavam alguma das cinco tecnologias medidas, ou 18\% quando ponderado por emprego; a adoção aparece agrupada com nuvem e robótica e concentrada em poucas cidades & Restrito aos Estados Unidos; os dados são de 2018 e antecedem a difusão da inteligência artificial generativa; adoção autorrelatada pela firma \\
\addlinespace
Ahmed e Wahed (2020) & Grandes empresas e universidades de elite ampliaram participação após 2012, deslocando universidades de médio e baixo escalão & Mede produção de pesquisa, não acesso comercial; o acesso a computação é inferido pela participação em conferências, não observado diretamente \emph{(inferida)} \\
\addlinespace
Sastry, Heim, Belfield \emph{et al.} (2024) & A computação é ponto de intervenção eficaz por ser detectável, excluível e quantificável, e por ser produzida em cadeia de suprimento concentrada & Preprint sem revisão por pares; propõe instrumentos de política sem avaliação empírica de efeito \emph{(inferida)} \\
\addlinespace
Mügge (2024) & A política prioriza independência jurisdicional sobre soberania do cidadão, adota a lógica de corrida tecnológica e negligencia efeitos sobre pessoas fora da União Europeia & Restrito a uma jurisdição; análise documental que não mede efeito sobre disponibilidade de serviço \emph{(inferida)} \\
\addlinespace
Furtado e Cunha (2024) & A expansão pressiona consumo de energia e água; a articulação entre discurso tecno-otimista e extração de recursos naturais pode aprofundar práticas coloniais & Não apresenta dado quantitativo de capacidade instalada; foco ambiental, sem tratar da disponibilidade de serviços \emph{(inferida)} \\
\addlinespace
\end{longtable}
\normalsize

\section{Síntese crítica}

\subsection{O que os trabalhos têm em comum}

Os dez trabalhos compartilham um pressuposto que nenhum deles enuncia explicitamente:
o de que o acesso desigual à tecnologia é um fenômeno \emph{mensurável}, e não apenas
uma condição a ser denunciada. Todos, inclusive os de natureza conceitual, organizam
sua argumentação em torno de uma quantidade --- nível de habilidade, densidade
tecnológica, taxa de adoção, participação em conferências, capacidade instalada --- que
poderia, em princípio, ser comparada entre unidades.

Há um segundo ponto comum, mais consequente. Todos localizam a desigualdade em uma
camada \emph{anterior} ao uso. Seja o equipamento \cite{vandeursen2018}, a habilidade
\cite{wang2024}, a infraestrutura nacional \cite{mannuru2023,furtado2024} ou o poder
computacional \cite{ahmed2020,sastry2024}, o argumento recorrente é que a diferença de
resultado observada decorre de uma diferença de dotação prévia. Essa convergência
importa porque define onde a medição deve incidir: não em quanto se usa, mas em quanto
está disponível para ser usado.

\subsection{Diferenças entre abordagens}

As diferenças, organizadas por eixo, são mais reveladoras do que as semelhanças.

\paragraph{Unidade de análise.} É o eixo de maior dispersão, e o que mais limita a
comparabilidade entre os trabalhos. \textcite{wang2024} e \textcite{lutz2019} operam
sobre o indivíduo; \textcite{mcelheran2024} e \textcite{ahmed2020} sobre a organização
--- firma e instituição de pesquisa; \textcite{knez2023} e \textcite{mannuru2023}
sobre o país; \textcite{mugge2024} sobre o bloco regulatório. Nenhum adota como unidade
o \emph{serviço em um território}, que é onde a disponibilidade de um modelo de IA
efetivamente se decide.

\begin{sloppypar}
\paragraph{Fonte de dados.} Quatro trabalhos são conceituais ou documentais
\cite{lutz2019,mannuru2023,sastry2024,mugge2024}; dois usam painéis macroeconômicos ou
censitários \cite{knez2023,mcelheran2024}; um usa survey \cite{wang2024}; um usa
bibliometria \cite{ahmed2020}; um usa corpus documental temático \cite{furtado2024}.
Apenas \textcite{amos2021} constrói sua própria série temporal a partir de documentos
arquivados na web. Essa distribuição revela uma dependência pesada de fontes que ou já
vêm agregadas por país, ou não são temporais.
\end{sloppypar}

\paragraph{Método.} O contraste mais nítido é entre a robustez causal e a proximidade
do objeto. \textcite{ahmed2020} aplicam controle sintético generalizado, o desenho
mais exigente do conjunto, mas sobre um objeto substituto: participação em
conferências como \emph{proxy} de acesso a computação. \textcite{mcelheran2024}
trabalham com o dado mais próximo da adoção real, mas apenas descritivamente e em um
único país. Nenhum trabalho combina desenho longitudinal com observação direta da
disponibilidade.

\subsection{Tendências observadas}

Três movimentos se desenham no conjunto.

O primeiro é o \textbf{deslocamento do objeto da desigualdade}. A literatura migrou do
acesso físico à rede para a qualidade do acesso material \cite{vandeursen2018}, daí
para habilidades e resultados \cite{lutz2019}, e mais recentemente para a capacidade
computacional \cite{ahmed2020,sastry2024}. Cada deslocamento acompanhou a mudança
daquilo que era escasso.

O segundo é a \textbf{entrada do Estado como ator de medição}. Nos trabalhos mais
recentes da terceira corrente \cite{sastry2024,mugge2024}, a desigualdade de acesso
deixa de ser efeito colateral do mercado e passa a ser instrumento e objeto de
política pública deliberada. Isso muda a natureza das variáveis explicativas
plausíveis: regime regulatório e alinhamento geopolítico tornam-se candidatos, ao lado
de renda e infraestrutura.

O terceiro é a \textbf{materialização do debate}. \textcite{furtado2024} recolocam a
infraestrutura física --- energia, água, território --- no centro de uma discussão que
vinha sendo conduzida em termos abstratos de soberania. A perspectiva é relevante
porque reintroduz o argumento concorrente mais forte: parte da assimetria observada
pode ser efeito de restrição material, e não de decisão política.

\subsection{Lacunas identificadas}

Do confronto entre as tabelas e os eixos acima resultam três lacunas.

\textbf{Primeira: a camada de computação de IA não é medida diretamente.} Os trabalhos
que tratam de poder computacional o fazem por substituição --- participação em
conferências \cite{ahmed2020} ou análise conceitual de instrumentos
\cite{sastry2024}. Os que medem adoção com dado empírico sólido
\cite{mcelheran2024,knez2023} não tratam de IA em nuvem. Não há, no conjunto revisado,
medição direta de \emph{quais modelos estão disponíveis em quais territórios e a
partir de quando}.

\textbf{Segunda: falta série temporal construída sobre documentação oficial.} A
disponibilidade de um serviço em uma região é informação publicada e datada, mas
volátil: a documentação é sobrescrita a cada atualização. \textcite{amos2021}
demonstram que arquivos da web podem ser convertidos em série longitudinal auditável
em escala de mais de um milhão de documentos, com validação e controle de qualidade.
O procedimento, contudo, foi aplicado a políticas de privacidade, e não a documentação
técnica de disponibilidade de serviço. A transposição não foi tentada.

\textbf{Terceira, e a mais consequente: nenhum trabalho separa atraso específico de
atraso geral.} Quando um serviço de IA demora a chegar a um território, a literatura
revisada não oferece instrumento para distinguir duas explicações concorrentes --- se o
atraso decorre da maturidade da infraestrutura local, caso em que todos os serviços da
mesma plataforma chegariam com atraso semelhante, ou se há tratamento diferenciado à
inteligência artificial. \textcite{knez2023} chega perto ao modelar custo relativo de
implementação, mas sobre tecnologias industriais e sem controle por serviço.
Sem essa separação, qualquer associação encontrada entre atraso e característica de
país permanece ambígua.

\subsection{Investigação futura}

As três lacunas convergem para um desenho de pesquisa viável.

Propõe-se tratar a defasagem como \textbf{tempo até evento}: a variável de interesse
seria o número de dias entre a primeira disponibilidade documentada de um modelo em
qualquer região e sua disponibilidade em uma região específica. A unidade de análise
passaria a ser o par formado por modelo e região --- unidade ausente de todos os
trabalhos revisados e que corresponde ao nível em que a decisão de disponibilização
efetivamente ocorre.

A formulação tem consequência metodológica direta. Regiões que nunca receberam
determinado modelo até o fim da janela de observação constituem observações
\emph{censuradas à direita}, e descartá-las enviesaria o resultado na direção de
subestimar a desigualdade, justamente nos casos extremos. O desenho apropriado é o de
análise de sobrevivência --- estimador de Kaplan-Meier estratificado, com modelo de
riscos proporcionais para o componente multivariado.

A fonte de dados seria a documentação técnica oficial das plataformas de nuvem,
recuperada retrospectivamente por meio de arquivos da web, aplicando à
documentação de disponibilidade o procedimento que \textcite{amos2021} aplicaram a
políticas de privacidade.

Por fim, a terceira lacuna exige uma variável de controle específica: o número total
de serviços da mesma plataforma disponíveis naquela região, como medida de maturidade
da infraestrutura local. Com ela, torna-se possível verificar se o atraso dos serviços
de inteligência artificial excede o atraso geral da região. Se exceder, há efeito
específico a explicar; se não exceder, o fenômeno é atribuível à maturidade da
infraestrutura. Essa separação é a contribuição que a literatura revisada não oferece
e é condição para que os resultados sejam interpretáveis.

\section{Link do projeto LaTeX}

O arquivo-fonte em \LaTeX{} utilizado na elaboração deste trabalho está disponível em:

\begin{center}
\large\texttt{https://www.overleaf.com/read/qbcvqyhkcymy}
\end{center}

\noindent O projeto está configurado com permissão de leitura para qualquer pessoa que
possua o endereço.

\printbibliography[title={Referências}]

\end{document}
